\section{建模提醒：不要只盯着神经网络外壳}
这节课反复强调 Bradley-Terry 和 MLE，其实是在提醒研究者：Reward Model 首先是一个概率模型，其次才是一个神经网络训练任务。只把注意力放在网络结构和 loss 调参上，很容易忽略偏好数据的统计假设是否成立。

\begin{warningbox}{Reward Model 常见误区}
\begin{itemize}
  \item 把偏好建模完全当成黑盒分类问题
  \item 不检查数据是否真的满足成对偏好假设
  \item 只看训练 loss，不分析 reward 分布是否稳定
\end{itemize}
\end{warningbox}

\subsection{本章小结}
掌握 Reward Model，关键是先理解偏好概率如何被建模，再决定用什么深度学习工具去拟合它。

\newpage
